\IfFileExists{stacks-project.cls}{%
\documentclass{stacks-project}
}{%
\documentclass{amsart}
}

% For dealing with references we use the comment environment
\usepackage{verbatim}
\newenvironment{reference}{\comment}{\endcomment}
%\newenvironment{reference}{}{}
\newenvironment{slogan}{\comment}{\endcomment}
\newenvironment{history}{\comment}{\endcomment}

% For commutative diagrams we use Xy-pic
\usepackage[all]{xy}

% We use 2cell for 2-commutative diagrams.
\xyoption{2cell}
\UseAllTwocells

% We use multicol for the list of chapters between chapters
\usepackage{multicol}

% This is generally recommended for better output
\usepackage{lmodern}
\usepackage[T1]{fontenc}

% For cross-file-references
\usepackage{xr-hyper}

% Package for hypertext links:
\usepackage{hyperref}

% For any local file, say "hello.tex" you want to link to please
% use \externaldocument[hello-]{hello}
\externaldocument[introduction-]{introduction}
\externaldocument[conventions-]{conventions}
\externaldocument[sets-]{sets}
\externaldocument[categories-]{categories}
\externaldocument[topology-]{topology}
\externaldocument[sheaves-]{sheaves}
\externaldocument[sites-]{sites}
\externaldocument[stacks-]{stacks}
\externaldocument[fields-]{fields}
\externaldocument[algebra-]{algebra}
\externaldocument[brauer-]{brauer}
\externaldocument[homology-]{homology}
\externaldocument[derived-]{derived}
\externaldocument[simplicial-]{simplicial}
\externaldocument[more-algebra-]{more-algebra}
\externaldocument[smoothing-]{smoothing}
\externaldocument[modules-]{modules}
\externaldocument[sites-modules-]{sites-modules}
\externaldocument[injectives-]{injectives}
\externaldocument[cohomology-]{cohomology}
\externaldocument[sites-cohomology-]{sites-cohomology}
\externaldocument[dga-]{dga}
\externaldocument[dpa-]{dpa}
\externaldocument[sdga-]{sdga}
\externaldocument[hypercovering-]{hypercovering}
\externaldocument[schemes-]{schemes}
\externaldocument[constructions-]{constructions}
\externaldocument[properties-]{properties}
\externaldocument[morphisms-]{morphisms}
\externaldocument[coherent-]{coherent}
\externaldocument[divisors-]{divisors}
\externaldocument[limits-]{limits}
\externaldocument[varieties-]{varieties}
\externaldocument[topologies-]{topologies}
\externaldocument[descent-]{descent}
\externaldocument[perfect-]{perfect}
\externaldocument[more-morphisms-]{more-morphisms}
\externaldocument[flat-]{flat}
\externaldocument[groupoids-]{groupoids}
\externaldocument[more-groupoids-]{more-groupoids}
\externaldocument[etale-]{etale}
\externaldocument[chow-]{chow}
\externaldocument[intersection-]{intersection}
\externaldocument[pic-]{pic}
\externaldocument[weil-]{weil}
\externaldocument[adequate-]{adequate}
\externaldocument[dualizing-]{dualizing}
\externaldocument[duality-]{duality}
\externaldocument[discriminant-]{discriminant}
\externaldocument[derham-]{derham}
\externaldocument[local-cohomology-]{local-cohomology}
\externaldocument[algebraization-]{algebraization}
\externaldocument[curves-]{curves}
\externaldocument[resolve-]{resolve}
\externaldocument[models-]{models}
\externaldocument[functors-]{functors}
\externaldocument[equiv-]{equiv}
\externaldocument[pione-]{pione}
\externaldocument[etale-cohomology-]{etale-cohomology}
\externaldocument[proetale-]{proetale}
\externaldocument[relative-cycles-]{relative-cycles}
\externaldocument[more-etale-]{more-etale}
\externaldocument[trace-]{trace}
\externaldocument[crystalline-]{crystalline}
\externaldocument[spaces-]{spaces}
\externaldocument[spaces-properties-]{spaces-properties}
\externaldocument[spaces-morphisms-]{spaces-morphisms}
\externaldocument[decent-spaces-]{decent-spaces}
\externaldocument[spaces-cohomology-]{spaces-cohomology}
\externaldocument[spaces-limits-]{spaces-limits}
\externaldocument[spaces-divisors-]{spaces-divisors}
\externaldocument[spaces-over-fields-]{spaces-over-fields}
\externaldocument[spaces-topologies-]{spaces-topologies}
\externaldocument[spaces-descent-]{spaces-descent}
\externaldocument[spaces-perfect-]{spaces-perfect}
\externaldocument[spaces-more-morphisms-]{spaces-more-morphisms}
\externaldocument[spaces-flat-]{spaces-flat}
\externaldocument[spaces-groupoids-]{spaces-groupoids}
\externaldocument[spaces-more-groupoids-]{spaces-more-groupoids}
\externaldocument[bootstrap-]{bootstrap}
\externaldocument[spaces-pushouts-]{spaces-pushouts}
\externaldocument[spaces-chow-]{spaces-chow}
\externaldocument[groupoids-quotients-]{groupoids-quotients}
\externaldocument[spaces-more-cohomology-]{spaces-more-cohomology}
\externaldocument[spaces-simplicial-]{spaces-simplicial}
\externaldocument[spaces-duality-]{spaces-duality}
\externaldocument[formal-spaces-]{formal-spaces}
\externaldocument[restricted-]{restricted}
\externaldocument[spaces-resolve-]{spaces-resolve}
\externaldocument[formal-defos-]{formal-defos}
\externaldocument[defos-]{defos}
\externaldocument[cotangent-]{cotangent}
\externaldocument[examples-defos-]{examples-defos}
\externaldocument[algebraic-]{algebraic}
\externaldocument[examples-stacks-]{examples-stacks}
\externaldocument[stacks-sheaves-]{stacks-sheaves}
\externaldocument[criteria-]{criteria}
\externaldocument[artin-]{artin}
\externaldocument[quot-]{quot}
\externaldocument[stacks-properties-]{stacks-properties}
\externaldocument[stacks-morphisms-]{stacks-morphisms}
\externaldocument[stacks-limits-]{stacks-limits}
\externaldocument[stacks-cohomology-]{stacks-cohomology}
\externaldocument[stacks-perfect-]{stacks-perfect}
\externaldocument[stacks-introduction-]{stacks-introduction}
\externaldocument[stacks-more-morphisms-]{stacks-more-morphisms}
\externaldocument[stacks-geometry-]{stacks-geometry}
\externaldocument[moduli-]{moduli}
\externaldocument[moduli-curves-]{moduli-curves}
\externaldocument[examples-]{examples}
\externaldocument[exercises-]{exercises}
\externaldocument[guide-]{guide}
\externaldocument[desirables-]{desirables}
\externaldocument[coding-]{coding}
\externaldocument[obsolete-]{obsolete}
\externaldocument[fdl-]{fdl}
\externaldocument[index-]{index}

% Theorem environments.
%
\theoremstyle{plain}
\newtheorem{theorem}[subsection]{Theorem}
\newtheorem{proposition}[subsection]{Proposition}
\newtheorem{lemma}[subsection]{Lemma}

\theoremstyle{definition}
\newtheorem{definition}[subsection]{Definition}
\newtheorem{example}[subsection]{Example}
\newtheorem{exercise}[subsection]{Exercise}
\newtheorem{situation}[subsection]{Situation}

\theoremstyle{remark}
\newtheorem{remark}[subsection]{Remark}
\newtheorem{remarks}[subsection]{Remarks}

\numberwithin{equation}{subsection}

% Macros
%
\def\lim{\mathop{\mathrm{lim}}\nolimits}
\def\colim{\mathop{\mathrm{colim}}\nolimits}
\def\Spec{\mathop{\mathrm{Spec}}}
\def\Hom{\mathop{\mathrm{Hom}}\nolimits}
\def\Ext{\mathop{\mathrm{Ext}}\nolimits}
\def\SheafHom{\mathop{\mathcal{H}\!\mathit{om}}\nolimits}
\def\SheafExt{\mathop{\mathcal{E}\!\mathit{xt}}\nolimits}
\def\Sch{\mathit{Sch}}
\def\Mor{\mathop{\mathrm{Mor}}\nolimits}
\def\Ob{\mathop{\mathrm{Ob}}\nolimits}
\def\Sh{\mathop{\mathit{Sh}}\nolimits}
\def\NL{\mathop{N\!L}\nolimits}
\def\CH{\mathop{\mathrm{CH}}\nolimits}
\def\proetale{{pro\text{-}\acute{e}tale}}
\def\etale{{\acute{e}tale}}
\def\QCoh{\mathit{QCoh}}
\def\Ker{\mathop{\mathrm{Ker}}}
\def\Im{\mathop{\mathrm{Im}}}
\def\Coker{\mathop{\mathrm{Coker}}}
\def\Coim{\mathop{\mathrm{Coim}}}

% Boxtimes
%
\DeclareMathSymbol{\boxtimes}{\mathbin}{AMSa}{"02}

%
% Macros for moduli stacks/spaces
%
\def\QCohstack{\mathcal{QC}\!\mathit{oh}}
\def\Cohstack{\mathcal{C}\!\mathit{oh}}
\def\Spacesstack{\mathcal{S}\!\mathit{paces}}
\def\Quotfunctor{\mathrm{Quot}}
\def\Hilbfunctor{\mathrm{Hilb}}
\def\Curvesstack{\mathcal{C}\!\mathit{urves}}
\def\Polarizedstack{\mathcal{P}\!\mathit{olarized}}
\def\Complexesstack{\mathcal{C}\!\mathit{omplexes}}
% \Pic is the operator that assigns to X its picard group, usage \Pic(X)
% \Picardstack_{X/B} denotes the Picard stack of X over B
% \Picardfunctor_{X/B} denotes the Picard functor of X over B
\def\Pic{\mathop{\mathrm{Pic}}\nolimits}
\def\Picardstack{\mathcal{P}\!\mathit{ic}}
\def\Picardfunctor{\mathrm{Pic}}
\def\Deformationcategory{\mathcal{D}\!\mathit{ef}}
\usepackage{cleveref}
\crefname{originalreference}{original reference}{original references}
\Crefname{originalreference}{Original reference}{Original references}
\makeatletter
\newlabel{bootstrap-theorem-final-bootstrap}{{04S6}{source}{}{sourcepointer-1}{}}
\newlabel{bootstrap-theorem-final-bootstrap@cref}{{[originalreference][0][]04S6}{[1][0][]source}{}{}{}}
\makeatother
\crefname{originallemma}{original lemma}{original lemmas}
\Crefname{originallemma}{Original lemma}{Original lemmas}
\makeatletter
\newlabel{lemma-base-change-hom-functor}{{05Y4}{source}{}{sourcepointer-2}{}}
\newlabel{lemma-base-change-hom-functor@cref}{{[originallemma][0][]05Y4}{[1][0][]source}{}{}{}}
\makeatother
\crefname{originallemma}{original lemma}{original lemmas}
\Crefname{originallemma}{Original lemma}{Original lemmas}
\makeatletter
\newlabel{lemma-etale-covering-hom-functor}{{05Y5}{source}{}{sourcepointer-3}{}}
\newlabel{lemma-etale-covering-hom-functor@cref}{{[originallemma][0][]05Y5}{[1][0][]source}{}{}{}}
\makeatother
\crefname{originalreference}{original reference}{original references}
\Crefname{originalreference}{Original reference}{Original references}
\makeatletter
\newlabel{more-morphisms-lemma-hom-from-finite-free-into-affine}{{05Y6}{source}{}{sourcepointer-4}{}}
\newlabel{more-morphisms-lemma-hom-from-finite-free-into-affine@cref}{{[originalreference][0][]05Y6}{[1][0][]source}{}{}{}}
\makeatother
\crefname{originalreference}{original reference}{original references}
\Crefname{originalreference}{Original reference}{Original references}
\makeatletter
\newlabel{spaces-lemma-base-change-representable-transformations-property}{{02WL}{source}{}{sourcepointer-5}{}}
\newlabel{spaces-lemma-base-change-representable-transformations-property@cref}{{[originalreference][0][]02WL}{[1][0][]source}{}{}{}}
\makeatother
\crefname{originalreference}{original reference}{original references}
\Crefname{originalreference}{Original reference}{Original references}
\makeatletter
\newlabel{spaces-lemma-glueing-algebraic-spaces}{{02WR}{source}{}{sourcepointer-6}{}}
\newlabel{spaces-lemma-glueing-algebraic-spaces@cref}{{[originalreference][0][]02WR}{[1][0][]source}{}{}{}}
\makeatother
\title{The relative Mor functor into any target is an algebraic space for finite locally free source}
\author{The Stacks Project Authors}
\date{stacks-project:criteria.tex:local-snapshot-20261012; GNU FDL 1.2 or later}
\begin{document}
\maketitle
\noindent\textit{Editorial note.} Footnote 'Some details omitted' covers the set-theoretic size verification that the coproduct of the F\_E is an algebraic space; retained verbatim and flagged for the reviewer.\par\medskip
\noindent\textit{Editorial note.} Steps the author leaves to the reader are flagged here for the reviewer: source line(s) 1709 carry an "omitted"/"follows easily"/"immediate from the definitions" phrase. Those sentences are retained verbatim; the displayed argument around them stays inside the counted span.\par\medskip
\noindent\textit{Allowed standard prerequisites (editorial).}\begin{itemize}
\item Base change of Mor. lemma-base-change-hom-functor (same chapter, named). Makes the etale reduction on B possible.
\item Etale cover of Mor. lemma-etale-covering-hom-functor (same chapter, named). Makes the etale reduction on X possible.
\item Mor into affine. more-morphisms-lemma-hom-from-finite-free-into-affine (cross-chapter, named). Provides algebraicity of each F\_E.
\item Glueing spaces. spaces-lemma-glueing-algebraic-spaces (cross-chapter, named). Assembles the open covering into an algebraic space.
\end{itemize}\medskip
\ifdefined\hypertarget\hypertarget{sourcepointer-1}{}\fi
\noindent\textit{Source reference.} Original reference 04S6: focused-seeds/geometry/S9-crit-hom-functor-algebraic-space/ref-tags.txt, source lines 5--5. Navigation pointer only: the counted proof cites this statement by \textbackslash{}ref as a standing input; the deduction the counted proof has to make is carried out inside the counted span, and no part of that deduction is delegated to this pointer.\par\medskip
\ifdefined\hypertarget\hypertarget{sourcepointer-2}{}\fi
\noindent\textit{Source reference.} Original lemma 05Y4: sources/stacks-criteria/pinned-original-source.tex, source lines 1590--1590. Navigation pointer only: the counted proof cites this statement by \textbackslash{}ref as a standing input; the deduction the counted proof has to make is carried out inside the counted span, and no part of that deduction is delegated to this pointer.\par\medskip
\ifdefined\hypertarget\hypertarget{sourcepointer-3}{}\fi
\noindent\textit{Source reference.} Original lemma 05Y5: sources/stacks-criteria/pinned-original-source.tex, source lines 1612--1612. Navigation pointer only: the counted proof cites this statement by \textbackslash{}ref as a standing input; the deduction the counted proof has to make is carried out inside the counted span, and no part of that deduction is delegated to this pointer.\par\medskip
\ifdefined\hypertarget\hypertarget{sourcepointer-4}{}\fi
\noindent\textit{Source reference.} Original reference 05Y6: focused-seeds/geometry/S9-crit-hom-functor-algebraic-space/ref-tags.txt, source lines 6--6. Navigation pointer only: the counted proof cites this statement by \textbackslash{}ref as a standing input; the deduction the counted proof has to make is carried out inside the counted span, and no part of that deduction is delegated to this pointer.\par\medskip
\ifdefined\hypertarget\hypertarget{sourcepointer-5}{}\fi
\noindent\textit{Source reference.} Original reference 02WL: focused-seeds/geometry/S9-crit-hom-functor-algebraic-space/ref-tags.txt, source lines 7--7. Navigation pointer only: the counted proof cites this statement by \textbackslash{}ref as a standing input; the deduction the counted proof has to make is carried out inside the counted span, and no part of that deduction is delegated to this pointer.\par\medskip
\ifdefined\hypertarget\hypertarget{sourcepointer-6}{}\fi
\noindent\textit{Source reference.} Original reference 02WR: focused-seeds/geometry/S9-crit-hom-functor-algebraic-space/ref-tags.txt, source lines 8--8. Navigation pointer only: the counted proof cites this statement by \textbackslash{}ref as a standing input; the deduction the counted proof has to make is carried out inside the counted span, and no part of that deduction is delegated to this pointer.\par\medskip
\begin{proposition}
\label{proposition-hom-functor-algebraic-space}
Let $S$ be a scheme. Let $Z \to B$ and $X \to B$ be morphisms of
algebraic spaces over $S$. If $Z \to B$ is finite locally free
then $\mathit{Mor}_B(Z, X)$ is an algebraic space.
\end{proposition}

\begin{proof}
Choose a scheme $B' = \coprod B'_i$ which is a disjoint union of
affine schemes $B'_i$ and an \'etale surjective morphism $B' \to B$.
We may also assume that $B'_i \times_B Z$ is the spectrum of a ring
which is finite free as a $\Gamma(B'_i, \mathcal{O}_{B'_i})$-module.
By
Lemma \ref{lemma-base-change-hom-functor}
and
Spaces, Lemma
\ref{spaces-lemma-base-change-representable-transformations-property}
the morphism $\mathit{Mor}_{B'}(Z', X') \to \mathit{Mor}_B(Z, X)$
is surjective \'etale. Hence by
Bootstrap, Theorem \ref{bootstrap-theorem-final-bootstrap}
it suffices to prove the proposition when $B = B'$ is a disjoint union of
affine schemes $B'_i$ so that each $B'_i \times_B Z$ is finite free
over $B'_i$. Then it actually suffices to prove the result for the restriction
to each $B'_i$. Thus we may assume that $B$ is affine and that
$\Gamma(Z, \mathcal{O}_Z)$ is a finite free $\Gamma(B, \mathcal{O}_B)$-module.

\medskip\noindent
Choose a scheme $X'$ which is a disjoint union of affine schemes and
a surjective \'etale morphism $X' \to X$. By
Lemma \ref{lemma-etale-covering-hom-functor}
the morphism $\mathit{Mor}_B(Z, X') \to \mathit{Mor}_B(Z, X)$
is representable by algebraic spaces, \'etale, and surjective.
Hence by
Bootstrap, Theorem \ref{bootstrap-theorem-final-bootstrap}
it suffices to prove the proposition when $X$ is a disjoint union
of affine schemes. This reduces us to the case discussed in the next
paragraph.

\medskip\noindent
Assume $X = \coprod_{i \in I} X_i$ is a disjoint union of affine
schemes, $B$ is affine, and that $\Gamma(Z, \mathcal{O}_Z)$ is a finite
free $\Gamma(B, \mathcal{O}_B)$-module. For any finite subset
$E \subset I$ set
$$
F_E = \mathit{Mor}_B(Z, \coprod\nolimits_{i \in E} X_i).
$$
By More on Morphisms,
Lemma \ref{more-morphisms-lemma-hom-from-finite-free-into-affine}
we see that $F_E$ is an algebraic space. Consider the morphism
$$
\coprod\nolimits_{E \subset I\text{ finite}} F_E
\longrightarrow
\mathit{Mor}_B(Z, X)
$$
Each of the morphisms
$F_E \to \mathit{Mor}_B(Z, X)$ is an open immersion, because it is
simply the locus parametrizing pairs $(a, b)$ where $b$ maps into
the open subscheme $\coprod\nolimits_{i \in E} X_i$ of $X$. Moreover,
if $T$ is quasi-compact, then for any pair $(a, b)$ the image
of $b$ is contained in $\coprod\nolimits_{i \in E} X_i$ for some
$E \subset I$ finite. Hence the displayed arrow is in fact an
open covering and we win\footnote{Modulo
some set theoretic arguments. Namely, we have to show that
$\coprod F_E$ is an algebraic space. This follows because
$|I| \leq \text{size}(X)$ and $\text{size}(F_E) \leq \text{size}(X)$
as follows from the explicit description of $F_E$ in the proof of
More on Morphisms,
Lemma \ref{more-morphisms-lemma-hom-from-finite-free-into-affine}.
Some details omitted.} by
Spaces, Lemma \ref{spaces-lemma-glueing-algebraic-spaces}.
\end{proof}

\end{document}
